\scanpage{44}
\margin{43}
\mspage{1}
\margin{Caution, notations in what follows relative to those of the preceding calculations, in $\alpha,\beta,\alpha',\beta',\zeta,\eta,\zeta'\ldots$}
Consider a conformal 2-sphere \inserted{$\Sigma$} with a three-element subset $J$,
\[
(1)\qquad J\subset\Sigma,\qquad J=(R_i)_{i\in J}
\]
so that the group $G_J$ acts on $\Sigma$ by automorphisms \deleted{oriented} preserving orientation, \deleted{and that} \deleted{$G_J\times\{1,\tau\}$ the group of} \inserted{whereas the group of automorphisms} of the situation $(\Sigma,J)$ that are conformal is
\[
(2)\qquad\Gamma_{\Sigma,J}=\Gamma\simeq G_J\times\{1,\tau\}
\]
\margin{\uline{NB} The antipodal map is \uline{not} part of $\Gamma$, since it does not stabilize $J$}
$\tau$ is the unique \deleted{anti} automorphism $\ne1$ \deleted{\unsure{which is}} which fixes the $R_i$ \deleted{\unsure{all}} ($\tau={}$complex conjugation).
\margin{thus the category of $(\Sigma,J)$ equivalent to that of pairs $(J,\uline{\bar\omega})$, where $J\in(\mathrm{Ens})_3$ and $\uline{\bar\omega}\in(\mathrm{Ens})_2$ cf below}
There is on $\Sigma$ a unique homogeneous Riemannian metric invariant under $G_J$, \inserted{up to a scalar $\lambda\in\mathbb R^{\unsure{*+|++}}$} \deleted{for which}, which allows us to identify $\Sigma$ canonically with the unit sphere in a real Euclidean vector space $E_3$. We put
\[
(3)\qquad\left\{\begin{aligned}
\uline{\bar\omega}&=\Omega(\Sigma)\simeq\Omega(E_3)\\
\uline{\omega}&=\Omega(J)\simeq\Omega(\Sigma_{\mathbb R})\\
\uline{\varepsilon}&=\Pi_0(\Sigma\setminus\Sigma_{\mathbb R})=\Pi_0(E\setminus V)=\Pi_0((E/V)^*)
\end{aligned}\right.
\]
\margin{set of the two orientations of $\Sigma$, or thus of $E$}
\margin{One could write $E(J,\uline{\bar\omega})$ $V(J,\uline{\omega})$.}
\margin{We have
\[
\uline{\omega}\wedge\uline{\varepsilon}\wedge\uline{\bar\omega}=1
\]
i.e.
\[
(3\text{ bis})\qquad\left\{\begin{aligned}
\uline{\omega}&\simeq\uline{\varepsilon}\wedge\uline{\bar\omega}\\
\uline{\varepsilon}&\simeq\uline{\bar\omega}\wedge\uline{\omega}\\
\uline{\bar\omega}&\simeq\uline{\omega}\wedge\uline{\varepsilon}
\end{aligned}\right.
\]
by the Stokes convention between orientations of a hemisphere and those of their boundary, more precise (4) here}
$V$ is the \inserted{$\mathbb R$-}vector subspace of $E$ \uline{spanned} by the $R_J$, which is of dimension 2, and where
\[
(4)\qquad\Sigma_{\mathbb R}=\Sigma\cap V
\]
Euclidean circle in the Euclidean plane $V$
